\section*{Ethics Statement}

This research adheres to established ethical standards for conducting work in high-stakes domains such as law. The legal documents used in our study were sourced from IndianKanoon (\url{https://indiankanoon.org/}), a publicly available repository of Indian court judgments. All documents are in the public domain and do not include sealed cases or personally identifiable sensitive information, ensuring that our use of the data complies with privacy and confidentiality norms.

We emphasize that the proposed \texttt{NyayaRAG} system is developed strictly for academic research purposes to simulate realistic legal reasoning processes. It is not intended for direct deployment in real-world legal settings. The model outputs must not be construed as legal advice, official court predictions, or determinants of legal outcomes. Any downstream use should be performed with oversight by qualified legal professionals. We strongly discourage the use of this system in live legal cases, policymaking, or decisions that may affect individuals' rights without appropriate human-in-the-loop supervision.

As part of our evaluation protocol, we involved domain experts (legal professionals and researchers) to assess the quality and legal coherence of the generated explanations. The evaluation was conducted on a curated subset of samples, and all participating experts were informed of the research objectives and voluntarily participated without any coercion or conflict of interest. No personal data was collected during this process, and all expert feedback was anonymized for analysis.

While we strive to enhance legal interpretability and transparency, we acknowledge that legal documents themselves may reflect systemic biases. Our framework, while replicating judicial reasoning patterns, may inherit such biases from training data. We do not deliberately introduce or amplify such biases, but we recognize the importance of further work in fairness auditing, particularly across litigant identity, socio-demographic markers, and jurisdictional diversity.